\documentclass[a4paper, serif, 10pt]{moderncv}

%% moderncv
% color and style
\moderncvcolor{black}
\moderncvstyle{banking}

%% page layout/margins
\usepackage[top=2.5cm, bottom=2.5cm, left=1.5cm, right=1.5cm]{geometry}

\nopagenumbers{}

%% Babel config for Italian language
% ref:
% - https://latex3.github.io/babel/guides/locale-italian.html
\usepackage[italian]{babel}

%% fontspec
\usepackage{fontspec}

\IfFontExistsTF{Linux Libertine}{
  \setmainfont[Ligatures={TeX, Common}, Numbers=OldStyle]{Linux Libertine}
}{}
\IfFontExistsTF{Linux Libertine O}{
  \setmainfont[Ligatures={TeX, Common}, Numbers=OldStyle]{Linux Libertine O}
}{}

%% CTeX
% CJK support, used to show CJK characters in the resume
%
% - heading=false: disable ctex heading
% - punct=kaiming: use kaiming punctuations styles for CJK
% - scheme=plain: use plain scheme, do not override \normalsize font size
% - space=auto: space settings for CJK characters
%
% CJK fonts are resolved per locale language, e.g. Noto Serif/Sans CJK JP
% for Japanese resumes, with fallback to the next region if not installed
%
% ref:
% - http://ctan.mirrorcatalogs.com/language/chinese/ctex/ctex.pdf
\usepackage[UTF8, heading=false, punct=kaiming, scheme=plain, space=auto]{ctex}

\IfFontExistsTF{Noto Serif CJK SC}{
  \setCJKmainfont{Noto Serif CJK SC}
}{}
\IfFontExistsTF{Noto Sans CJK SC}{
  \setCJKsansfont{Noto Sans CJK SC}
}{}

%% line spacing
\usepackage{setspace}
\setstretch{1.125}

%% URL styling - use normal text instead of monospace
\urlstyle{same}

%% Patch \httplink and \httpslink to handle full URLs
% The original moderncv macros always prepend http:// or https:// to URLs.
% This causes issues when the URL already has a protocol (e.g., https://example.com)
% resulting in malformed URLs like https://https://example.com.
% This patch checks if the URL already contains "://" and uses it directly if so.
% Note: We use \str_set:Nx (x-expansion) to expand macros like \@homepage first.
\ExplSyntaxOn
\str_new:N \g_yamlresume_url_str

\RenewDocumentCommand{\httpslink}{O{}m}{
  \str_set:Nx \g_yamlresume_url_str {#2}
  \str_if_in:NnTF \g_yamlresume_url_str {://}
    { \url{#2} }
    { \url{https://#2} }
}

\RenewDocumentCommand{\httplink}{O{}m}{
  \str_set:Nx \g_yamlresume_url_str {#2}
  \str_if_in:NnTF \g_yamlresume_url_str {://}
    { \url{#2} }
    { \url{http://#2} }
}
\ExplSyntaxOff

\name{Luca Bianchi}{}
\title{QA Engineer}
\phone[mobile]{+39 333 123 4567}
\email{luca.bianchi@example.com}
\homepage{https://www.lucabianchi.it}

\address{Via Roma 1, Milano, Lombardia, Italia, 20100}{}{}

\extrainfo{{\small \faLinkedin}\ \href{https://www.linkedin.com/in/lucabianchi}{@lucabianchi} {} {} {} • {} {} {} 
{\small \faGithub}\ \href{https://github.com/lucabianchi}{@lucabianchi}}

\begin{document}

\maketitle

\section{Informazioni personali}

\cvline{}{\begin{itemize}
\item QA Engineer con oltre 5 anni di esperienza nella verifica e validazione di software
\item Specializzato in test manuali e automatizzati, con solida conoscenza di metodologie Agile
\item Esperto nell'utilizzo di strumenti come Selenium, JUnit, TestNG e Jenkins
\item Capacità di analisi dei requisiti e scrittura di test case dettagliati
\item Forte attenzione ai dettagli e orientamento alla qualità del prodotto
\end{itemize}}
\section{Formazione}

\cventry{set 2018–lug 2020}
        {Laurea magistrale, Ingegneria Informatica, Punteggio: 110/110}
        {Politecnico di Milano}
        {\url{https://www.polimi.it}}
        {}
        {\begin{itemize}
\item Approfondimento delle metodologie di testing e quality assurance
\item Partecipazione a progetti di gruppo per lo sviluppo di software di qualità
\item Tesi sperimentale su tecniche di test automatico per applicazioni web
\end{itemize}
\textbf{Corsi}: Ingegneria del Software, Sistemi di Qualità, Metodi Formali per il Software, Testing e Verifica del Software, Architetture Software, Sicurezza Informatica}
\section{Esperienza professionale}

\cventry{gen 2021–Presente}
        {Senior QA Engineer}
        {TechQuality S.p.A.}
        {\url{https://www.techquality.it}}
        {}
        {\begin{itemize}
\item Responsabile della strategia di test per un'importante piattaforma e-commerce
\item Progettazione e implementazione di suite di test automatici con Selenium e TestNG
\item Coordinamento di un team di 3 QA Engineer junior
\item Collaborazione con i team di sviluppo per l'analisi dei difetti e il miglioramento continuo
\item Riduzione del 30\% dei difetti in produzione grazie a test di regressione automatizzati
\end{itemize}
\textbf{Parole chiave}: Automazione dei test, Selenium, TestNG, Jenkins, Agile, Qualità del software}

\cventry{set 2018–dic 2020}
        {QA Engineer}
        {SoftwareTest Italia}
        {\url{https://www.softwaretestitalia.it}}
        {}
        {\begin{itemize}
\item Esecuzione di test manuali e automatici su applicazioni web e mobile
\item Scrittura di test case e piani di test basati sui requisiti funzionali
\item Utilizzo di JIRA per la gestione dei difetti e la reportistica
\item Partecipazione a sprint review e daily stand-up in ambiente Agile
\item Collaborazione con sviluppatori per la risoluzione di bug critici
\end{itemize}
\textbf{Parole chiave}: Test manuale, Test automatico, JIRA, Selenium, Postman, Agile}
\section{Lingue}

\cvline{Italiano}{Madrelingua o bilingue \hfill \textbf{Parole chiave}: Madrelingua}
\cvline{Inglese}{Competenza professionale completa \hfill \textbf{Parole chiave}: IELTS 7.5, Inglese commerciale}
\section{Competenze}

\cvline{Test Automation}{Esperto \hfill \textbf{Parole chiave}: Selenium, TestNG, JUnit, Cucumber, Jenkins, Postman}
\cvline{Test Manuale}{Avanzato \hfill \textbf{Parole chiave}: Analisi dei requisiti, Scrittura di test case, Test di regressione, Test di integrazione, Test di sistema}
\cvline{Metodologie}{Avanzato \hfill \textbf{Parole chiave}: Agile, Scrum, Kanban, TDD, BDD}
\cvline{Strumenti}{Avanzato \hfill \textbf{Parole chiave}: JIRA, Confluence, Git, SQL, Docker}
\section{Premi}

\cventry{dic 2022}
        {TechQuality S.p.A.}
        {Premio Qualità 2022}
        {}
        {}
        {Riconoscimento per l'eccellenza nella implementazione di processi di test automatizzati che hanno portato a una riduzione significativa dei difetti in produzione.}
\section{Certificazioni}

\cventry{mar 2021}
        {ISTQB}
        {ISTQB Certified Tester - Advanced Level}
        {\url{https://www.istqb.org/certifications/advanced-level}}
        {}
        {}

\cventry{giu 2020}
        {Selenium Academy}
        {Certified Selenium Professional}
        {\url{https://www.seleniumacademy.com/certification}}
        {}
        {}
\section{Pubblicazioni}

\cventry{set 2021}
        {Automazione dei test per applicazioni web moderne}
        {Rivista Italiana di Qualità del Software}
        {\url{https://www.qa-rivista.it/articoli/automazione-test-web}}
        {}
        {\begin{itemize}
\item Articolo che esplora le migliori pratiche per l'automazione dei test in ambienti Agile
\item Discussione di casi di studio e strumenti consigliati
\end{itemize}}
\section{Progetti}

\cventry{gen 2022–giu 2022}
        {Framework di test automatizzato per applicazioni web basato su Selenium e TestNG}
        {Framework di Automazione dei Test}
        {\url{https://github.com/lucabianchi/test-framework}}
        {}
        {\begin{itemize}
\item Sviluppo di un framework modulare per test automatici riutilizzabili
\item Integrazione con Jenkins per l'esecuzione continua dei test
\item Riduzione del tempo di esecuzione dei test del 40\%
\end{itemize}
\textbf{Parole chiave}: Selenium, TestNG, Jenkins, Automazione}

\cventry{mar 2021–ago 2021}
        {Dashboard per il monitoraggio delle metriche di qualità del software}
        {Dashboard per le Metriche di Qualità}
        {}
        {}
        {\begin{itemize}
\item Creazione di una dashboard interattiva per visualizzare i risultati dei test
\item Integrazione con JIRA per il recupero automatico dei dati sui difetti
\item Miglioramento della visibilità sulla qualità del prodotto per il team
\end{itemize}
\textbf{Parole chiave}: Dashboard, Metriche, JIRA, Qualità}
\section{Interessi}

\cvline{Sport}{Corsa, Nuoto, Escursionismo}
\cvline{Tecnologia}{Open Source, Intelligenza Artificiale, Scrittura di blog}
\section{Volontariato}

\cventry{set 2021–Presente}
        {Mentore QA}
        {Associazione Qualità Italia}
        {\url{https://www.qualityitalia.it}}
        {}
        {\begin{itemize}
\item Volontariato come mentor per giovani QA Engineer
\item Organizzazione di workshop su test automation e quality assurance
\item Supporto nella preparazione alle certificazioni ISTQB
\end{itemize}}
    
\end{document}